\newpage
\section{Misión 6: Molestar a dos planetas enteros}\label{sec:mission6}

\begin{minipage}{0.87\linewidth}
Han pasado años en la colonia desde que ayudaste a configurar la gran antena que os permite comunicaros con la Tierra. Hace poco, tu colonia ha descubierto \textit{otra} colonia secreta, esta en Saturno -- ¡y la Tierra ni se molestó en decírnoslo!\\

\noindent
La asamblea ha decidido gastar una pequeña broma y te pide ayuda de nuevo.
Tu objetivo es secuestrar la próxima conexión TCP Saturno $\leftrightarrow$ Tierra
y enviar un comunicado falso haciéndote pasar por Saturno. La asamblea pide que el comunicado tenga clase.
\end{minipage}
%
\hspace{0.03\linewidth}
%
\begin{minipage}{0.10\linewidth}
\includegraphics[width=\linewidth]{ostrich}
\end{minipage}

\subsection*{Detalles del informe}
\noindent
La Ingeniera Jefe de la colonia ha hecho reconocimiento y está dispuesta a ayudarte.
Ha reunido los siguientes datos:

\begin{itemize}[itemsep=0.1cm]
\item Tu HugeAntenna\texttrademark\ de la colonia se comunica con otra HugeAntenna\texttrademark\ de la Tierra. La Tierra usa esa misma antena para comunicarse con otra HugeAntenna\texttrademark\ de Saturno.
Cuando la Tierra recibe el primer paquete de una conexión TCP, verifica su origen.

\item Tras ese primer paquete, la Tierra no verifica el origen de ningún paquete. De hecho, los mensajes falsos o inesperados recibidos después se ignoran sin más, sin activar ninguna alarma. Con sincronización cuidadosa, usarás esto para inyectar tus paquetes falsificados.

\item El primer día de cada mes, a las 4:20 a.\,m. en punto (hora de Nueva Barcelona, Tierra), la Tierra recibe el primer paquete de una conexión TCP de Saturno. Estas conexiones sirven para intercambiar comunicados oficiales, y la asamblea quiere que secuestres la próxima.

\item Saturno usa una pila TCP/IP anticuada e insegura con este fallo: siempre fija a $0$ el número de secuencia de TCP en el primer mensaje de una conexión TCP.

\item No podrás ver los mensajes intercambiados entre la Tierra y Saturno. Sin embargo, la Ingeniera Jefe logró capturar la siguiente frame (trama) Ethernet, que inició la conexión TCP anterior (llamémosla $C_{-1}$) entre Saturno y la Tierra (\texttt{00 d0} va seguido de \texttt{f1 bb} y \texttt{45 10} va seguido de \texttt{00 28}):

\vspace{0.25cm}
\begin{center}
\begin{minipage}{0.5\linewidth}
\begin{verbatim}
00 09 bf 42 07 64 00 d0   f1 bb 0f 5a 08 00 45 10
00 28 d2 71 40 00 40 06   00 00 ca f7 44 05 4c 4c
15 15 9f ee 00 a9 00 00   00 00 00 00 00 00 50 02
01 00 00 00 00 00
\end{verbatim}
\end{minipage}
\end{center}
\vspace{0.25cm}
\end{itemize}

\noindent
Con esto en mente, la asamblea solicita un informe que contenga los siguientes aspectos:

\begin{enumerate}[itemsep=0.1cm]
\item Dibuja un diagrama de carriles del establecimiento de la conexión anterior ($C_{-1}$) entre Saturno y la Tierra.
Además, tabula \textit{todos} los campos de cabecera presentes en estos paquetes (y sus valores).

\item Explica tu plan de ataque, \textit{i.e.}, cómo piensas secuestrar la conexión. Indica claramente:
    \begin{itemize}
    \item Qué suposiciones haces sobre la próxima conexión (p. ej., qué campos seguirán idénticos o podrían cambiar).
    \item Cuántos mensajes enviarás y su sincronización prevista.
    \end{itemize}

\item Muestra la secuencia completa de bytes de una frame que contenga tu comunicado de broma.
\end{enumerate}

\subsection*{Detalles del feedback}

Emplea un proceso de feedback (retroalimentación) similar al de las misiones anteriores, abordando estos aspectos clave:
\begin{itemize}
\item ¿Se decodificó correctamente la frame capturada por la Ingeniera Jefe?
\item ¿Lograste idear un plan que requiera menos mensajes?
\item ¿Pudiste decodificar el comunicado de broma? ¿Tenía clase?
\end{itemize}

